\section{Objetivo General}
Desarrollar la reingeniería mecánica y la automatización de un brazo robótico cartesiano mediante el rediseño de sus sistemas de traslación y la integración de visión computacional, sensores y control embebido, con el fin de ejecutar tareas de detección, manipulación, clasificación y transferencia de productos en la línea de empaque del Laboratorio de Diseño de Procesos de la Universidad del Valle de Guatemala.

\section{Objetivos Específicos}
\begin{enumerate}
\item Rediseñar los componentes estructurales y los sistemas de traslación de los ejes X, Y, Z mediante la identificación de los puntos críticos de fricción y el uso de software de diseño Autodesk Inventor, con el fin de integrar rodamientos lineales y ejes endurecidos que aseguren la fluidez cinemática. 
\item Desarrollar un módulo de operación manual mediante la programación de funciones de movimiento individual por eje y el uso de un mando inalámbrico, para permitir el posicionamiento preventivo y la gestión de mantenimiento.
\item Desarrollar un sistema de control de lazo cerrado mediante la integración de un sistema de visión, sensores de posición y un codificador rotativo en la banda transportadora, para automatizar la detección de objetos y sincronizar el movimiento del brazo con la velocidad de producción.
\item Evaluar la precisión y repetibilidad del brazo robótico, con el fin de validar las mejoras mecánicas y el sistema de control mediante pruebas repetitivas. 
\end{enumerate}